\title{Lecture 14\\Intelligent Systems Design Technologies \vspace{-2em}}
%\author[]{Shunkevich D.V.}
%\institute[]{Belarusian State University of Informatics and Radioelectronics}

\begin{frame}
	\titlepage
\end{frame}

\begin{frame}{\\Lecture Contents}
	\topline
	\justifying
	\begin{SCn}
		
		Lecture 14. Intelligent Systems Design Technologies
		\begin{scnrelfromset}{structure}
			\scnitem{Reusable components of an intelligent system}
			\scnitem{Technology for designing intelligent computer systems}
			\scnitem{Requirements for the design of intelligent systems}
		\end{scnrelfromset}
		
	\end{SCn}
\end{frame}

\begin{frame}{Reusable components of an intelligent system}
	\topline
	\justifying
	\begin{SCn}
		
		\scnheader{reusable components in intelligent systems}
		\scnidtf{components that can be reused across different systems or across different projects within the same system}
		
	\end{SCn}
\end{frame}

\begin{frame}{Reusable components of an intelligent system}
	\topline
	\justifying
	\begin{SCn}
		
		\scnheader{reusable components in intelligent systems}
		
		\begin{scnrelfromset}{application features}
			\scnfileitem{Using it significantly reduces the time and cost of developing new systems or components.}
			\scnfileitem{Reusable components are used for data processing, model tuning, feature extraction, and other tasks. Components are developed and tested individually and can then be incorporated into various machine learning systems without additional work. This helps reduce development time and code reuse.}
		\end{scnrelfromset}
		
	\end{SCn}
\end{frame}

\begin{frame}{Reusable components of an intelligent system}
	\topline
	\justifying
	\begin{SCn}
		\scnheader{reusable components in intelligent systems}
		
		\begin{scnrelfromset}{application features}
			\scnfileitem{Used to create a common knowledge base that can be used by many systems}
			\scnfileitem{Reusable components are used to unify interfaces and standards between different systems, increasing compatibility and data exchange. For example, speech recognition systems feature reusable components for audio processing and speech feature extraction. These components can be shared across multiple systems to make speech processing more unified and user-friendly.}
		\end{scnrelfromset}
		
	\end{SCn}
\end{frame}

\begin{frame}{Reusable components of an intelligent system}
	\topline
	\justifying
	\begin{SCn}
		
		Component-Based Design of Intelligent Systems
		\begin{scnrelfromset}{requirements}
			\scnfileitem{Using a universal knowledge representation language used in an intelligent system.}
			\scnfileitem{The presence of a universal procedure for integrating knowledge within the specified language.}
			\scnfileitem{The presence of a standard that ensures the semantic compatibility of integrated knowledge (such a standard is a consistent system of used concepts and a unified specification of components).}
		\end{scnrelfromset}
		
	\end{SCn}
\end{frame}

\begin{frame}{Reusable components of an intelligent system}
	\topline
	\justifying
	\begin{SCn}
		
		Component-Based Design of Intelligent Systems
		\begin{scnrelfromset}{requirements}
			\scnfileitem{A reusable component should perform its functions in the most general way possible so that the range of possible systems into which it can be embedded is as wide as possible.}
			\scnfileitem{Self-sufficiency of components (or groups of components) of a technology, i.e. their ability to function separately from other components without losing the feasibility of their use.}
		\end{scnrelfromset}
		
	\end{SCn}
\end{frame}

\begin{frame}{Technology for designing intelligent computer systems}
	\topline
	\justifying
	\begin{SCn}
		
		\scnheader{Technology for designing intelligent computer systems}
		\scnidtf{a set of methods and tools used to create systems capable of self-learning, data analysis, and decision making based on available information}
		\scntext{examples}{natural language processing systems, machine learning systems, computer vision systems, expert systems, robotics systems.}
		
	\end{SCn}
\end{frame}

\begin{frame}{Technology for designing intelligent computer systems}
	\topline
	\justifying
	\begin{SCn}
		
		\scnheader{Technology for designing intelligent computer systems}
		\begin{scnrelfromset}{key aspects}
			\scnfileitem{The ability to analyze large volumes of data and use it for decision making. This technology can utilize various machine learning techniques, such as classification, clustering, regression, and natural language processing, which facilitate the training of computer systems based on large volumes of data.}
		\end{scnrelfromset}
		
	\end{SCn}
\end{frame}

\begin{frame}{Technology for designing intelligent computer systems}
	\topline
	\justifying
	\begin{SCn}
		
		\scnheader{Technology for designing intelligent computer systems}
		\begin{scnrelfromset}{key aspects}
			\scnfileitem{the ability to create systems that can learn from experience in real time. Such systems can be configured to behave in a specific way and can change their behavior based on new experience gained during operation.}
			\scnfileitem{includes various methods and techniques for creating intuitive interfaces for interacting with computer systems. These intuitive interfaces make using systems more convenient and understandable for people}
		\end{scnrelfromset}
		
	\end{SCn}
\end{frame}

\begin{frame}{Requirements for intelligent systems design technology}
	\topline
	\justifying
	\begin{SCn}
		
		\scnheader{Technology for designing intelligent computer systems}
		\begin{scnrelfromset}{requirements}
			\scnfileitem{Implementation of the proposed technology for the development and maintenance of next-generation intelligent computer systems in the form of an intelligent computer metasystem that fully complies with the standards of the proposed next-generation intelligent computer systems developed using the proposed technology}
			\scnfileitem{Unification and standardization of new-generation intelligent computer systems, as well as methods for their design, implementation, maintenance, reengineering, and operation}
		\end{scnrelfromset}
		
	\end{SCn}
\end{frame}

\begin{frame}{Requirements for intelligent systems design technology}
	\topline
	\justifying
	\begin{SCn}
		
		\scnheader{Technology for designing intelligent computer systems}
		\begin{scnrelfromset}{requirements}
			\scnfileitem{Component-based design of next-generation intelligent computer systems, i.e., design focused on assembling intelligent computer systems from ready-made components based on continuously expanding libraries of reusable components.}
			\scnfileitem{Permanent evolution of the standard for next-generation intelligent computer systems, as well as methods for their design, implementation, maintenance, reengineering, and operation}
		\end{scnrelfromset}
		
	\end{SCn}
\end{frame}

\begin{frame}{Requirements for intelligent systems design technology}
	\topline
	\justifying
	\begin{SCn}
		
		\scnheader{Technology for designing intelligent computer systems}
		\begin{scnrelfromset}{requirements}
			\scnfileitem{Clear coordination and prompt formalized recording (in the form of formal ontologies) of the approved current state of the hierarchical system of all concepts underlying the permanently evolving standard of next-generation intelligent computer systems, as well as the basis of each intelligent computer system being developed}
			\scnfileitem{A sufficiently complete and timely documentation of the current state of each project}
			\scnfileitem{Using a top-down design approach{}}
		\end{scnrelfromset}
		
	\end{SCn}
\end{frame}

\begin{frame}{Requirements for intelligent systems design technology}
	\topline
	\justifying
	\begin{SCn}
		
		\scnheader{Technology for designing intelligent computer systems}
		\begin{scnrelfromset}{requirements}
			\scnfileitem{The complex nature of the proposed technology, which supports the design of not only components of new-generation intelligent computer systems (various fragments of knowledge bases, knowledge bases in general, various methods for solving problems, various internal information agents, problem solvers in general, interfaces in general), but also intelligent computer systems as a whole as independent design objects, taking into account the specifics of the classes to which the designed intelligent computer systems belong}
		\end{scnrelfromset}
		
	\end{SCn}
\end{frame}

\begin{frame}{Requirements for intelligent systems design technology}
	\topline
	\justifying
	\begin{SCn}
		
		\scnheader{Technology for designing intelligent computer systems}
		\begin{scnrelfromset}{requirements}
			\scnfileitem{The comprehensive nature of the proposed technology, which supports not only the integrated design of next-generation intelligent computer systems, but also their implementation (assembly, reproduction), maintenance, reengineering during operation, and the operation itself}
		\end{scnrelfromset}
		
	\end{SCn}
\end{frame}
